\section{Discussion}

The simulation illustrates why a successful rerun in the original environment
is an incomplete reproducibility test. Pipelines with fully captured environments
were substantially more resilient to dependency and system perturbations, and
explicit seed control improved numerical agreement within each environment
strategy. A compact evaluation should therefore report at least three elements:
execution success, numerical agreement, and the conditions under which both were
measured.

These findings align with calls to treat openness and reproducibility as
observable practices rather than aspirations \citep{nosek2015}. The proposed
framework adds a stress-testing layer: authors declare expected outputs, rebuild
the analysis under controlled changes, and report both failures and deviations.
This approach can be integrated into continuous-integration systems and repeated
when dependencies, data, or analysis code change.

\subsection{Limitations}

The study uses synthetic projects and stylized perturbations, so its numerical
estimates should not be interpreted as prevalence estimates for published
research. The simulated pipelines were also smaller than many production
workflows and did not represent specialized hardware, restricted data, or
interactive software. Empirical validation should sample real projects across
disciplines, define field-specific tolerances, and distinguish failures caused by
missing infrastructure from those caused by insufficient documentation.

\subsection{Implications for authors and reviewers}

For authors, the main practical implication is to preserve both the environment
and the sources of stochastic variation. For reviewers, a failed rebuild is most
informative when the submission records the attempted environment and the first
divergent stage. For journals, a short machine-readable manifest could accompany
the conventional data and code availability statements without imposing a
single software platform.
